% ============================================================
% title.tex - Титульный лист
% ============================================================
% !TeX program = xelatex

\begin{titlepage}
	\centering
	
	{\scshape\Large [Название учебного заведения / института]\par}
	\vspace{0.5cm}
	{\scshape\large [Название факультета или кафедры]\par}
	
	\vspace{3cm}
	
	\rule{\textwidth}{1.6pt}\vspace*{-\baselineskip}\vspace*{2pt}
	\rule{\textwidth}{0.4pt}\\[\baselineskip]
	
	{\LARGE \textbf{Методическое пособие}}
	
	\vspace{0.5cm}
	{\huge \textbf{Обработка данных ТГц-спектроскопии на языке Python}}
	
	\vspace{0.3cm}
	{\Large (на примере анализа проволочных поляризаторов)}
	
	\rule{\textwidth}{0.4pt}\vspace*{-\baselineskip}\vspace{3.2pt}
	\rule{\textwidth}{1.6pt}\\[\baselineskip]
	
	\vspace{2cm}
	
	\begin{flushright}
		{\large \textbf{Разработчик:}\\[0.2cm] [Ваши ФИО / Должность]}
	\end{flushright}
	
	\vspace{2cm}
	
	\begin{center}
		\fbox{\parbox{0.85\textwidth}{
				\centering
				\vspace{0.2cm}
				\textbf{\large Аннотация}
				
				\vspace{0.3cm}
				Данное пособие представляет собой пошаговое руководство по созданию
				программы для обработки данных терагерцовой спектроскопии.
				Материал ориентирован на студентов инженерных и физических специальностей,
				желающих освоить Python для решения реальных прикладных задач.
				В пособии применяется подход «от простого к сложному» и разбирается 
				полный цикл цифровой обработки сигнала: от чтения файлов до проверки 
				физических моделей.
				\vspace{0.2cm}
		}}
	\end{center}
	
	\vfill
	
	{\large \textbf{Ключевые слова:} Python, NumPy, SciPy, Matplotlib, ТГц-спектроскопия, цифровая обработка сигналов, БПФ, проволочные поляризаторы, модель Бланко.}
	
	\vspace{1cm}
	{\large \textbf{\the\year{}} г.}
	
\end{titlepage}
